\documentclass[11pt]{article}

% Change "review" to "final" to generate the final (sometimes called camera-ready) version.
% Change to "preprint" to generate a non-anonymous version with page numbers.
\usepackage[review]{acl}

% Standard package includes
\usepackage{times}
\usepackage{latexsym}

% For proper rendering and hyphenation of words containing Latin characters (including in bib files)
\usepackage[T1]{fontenc}
% For Vietnamese characters
% \usepackage[T5]{fontenc}
% See https://www.latex-project.org/help/documentation/encguide.pdf for other character sets

% This assumes your files are encoded as UTF8
\usepackage[utf8]{inputenc}

% This is not strictly necessary, and may be commented out,
% but it will improve the layout of the manuscript,
% and will typically save some space.
\usepackage{microtype}

% This is also not strictly necessary, and may be commented out.
% However, it will improve the aesthetics of text in
% the typewriter font.
\usepackage{inconsolata}

%Including images in your LaTeX document requires adding
%additional package(s)
\usepackage{graphicx}

% If the title and author information does not fit in the area allocated, uncomment the following
%
%\setlength\titlebox{<dim>}
%
% and set <dim> to something 5cm or larger.

\title{Commit, Then Merge: Subgroup-Conditioned Drafts\\for Overton-Pluralistic Answers}

% Author information can be set in various styles:
% For several authors from the same institution:
% \author{Author 1 \and ... \and Author n \\
%         Address line \\ ... \\ Address line}
% if the names do not fit well on one line use
%         Author 1 \\ {\bf Author 2} \\ ... \\ {\bf Author n} \\
% For authors from different institutions:
% \author{Author 1 \\ Address line \\  ... \\ Address line
%         \And  ... \And
%         Author n \\ Address line \\ ... \\ Address line}
% To start a separate ``row'' of authors use \AND, as in
% \author{Author 1 \\ Address line \\  ... \\ Address line
%         \AND
%         Author 2 \\ Address line \\ ... \\ Address line \And
%         Author 3 \\ Address line \\ ... \\ Address line}

\author{First Author \\
  Affiliation / Address line 1 \\
  Affiliation / Address line 2 \\
  Affiliation / Address line 3 \\
  \texttt{email@domain} \\\And
  Second Author \\
  Affiliation / Address line 1 \\
  Affiliation / Address line 2 \\
  Affiliation / Address line 3 \\
  \texttt{email@domain} \\}

%\author{
%  \textbf{First Author\textsuperscript{1}},
%  \textbf{Second Author\textsuperscript{1,2}},
%  \textbf{Third T. Author\textsuperscript{1}},
%  \textbf{Fourth Author\textsuperscript{1}},
%\\
%  \textbf{Fifth Author\textsuperscript{1,2}},
%  \textbf{Sixth Author\textsuperscript{1}},
%  \textbf{Seventh Author\textsuperscript{1}},
%  \textbf{Eighth Author \textsuperscript{1,2,3,4}},
%\\
%  \textbf{Ninth Author\textsuperscript{1}},
%  \textbf{Tenth Author\textsuperscript{1}},
%  \textbf{Eleventh E. Author\textsuperscript{1,2,3,4,5}},
%  \textbf{Twelfth Author\textsuperscript{1}},
%\\
%  \textbf{Thirteenth Author\textsuperscript{3}},
%  \textbf{Fourteenth F. Author\textsuperscript{2,4}},
%  \textbf{Fifteenth Author\textsuperscript{1}},
%  \textbf{Sixteenth Author\textsuperscript{1}},
%\\
%  \textbf{Seventeenth S. Author\textsuperscript{4,5}},
%  \textbf{Eighteenth Author\textsuperscript{3,4}},
%  \textbf{Nineteenth N. Author\textsuperscript{2,5}},
%  \textbf{Twentieth Author\textsuperscript{1}}
%\\
%\\
%  \textsuperscript{1}Affiliation 1,
%  \textsuperscript{2}Affiliation 2,
%  \textsuperscript{3}Affiliation 3,
%  \textsuperscript{4}Affiliation 4,
%  \textsuperscript{5}Affiliation 5
%\\
%  \small{
%    \textbf{Correspondence:} \href{mailto:email@domain}{email@domain}
%  }
%}

\begin{document}
\maketitle
\begin{abstract}
A single answer to a contested question tends to articulate the majority
position and gesture at the rest. We study two ways of feeding retrieved survey
data into an answer built from several drafts and a lossless merge. In the
first, every draft sees all retrieved viewpoints and the model decides what to
foreground. In the second, a fixed rule assigns each draft a single survey
subgroup, and the draft argues that subgroup's view without hedging. Each is
compared against the same drafting and merging with no retrieved data, so that
any gain can be credited to the survey data rather than to drafting and merging
themselves. We evaluate by how many viewpoint clusters a single answer properly
represents.
\end{abstract}

\section{Method}

\subsection{Viewpoint graph}

We build the graph from population survey data, in our experiments the Pew
American Trends Panel: 1{,}492 questions, each answered by 58 demographic
subpopulations. The graph is a tree. Its leaves are \emph{opinion leaves}, one
for each subgroup's answers to one question; a leaf records, for instance, how
self-described conservatives answered a question on trust in government, as a
distribution $p_\ell$ over that question's answer options. Leaves sit under the
demographic axis they belong to (party, age, religion, and so on), axes under
their survey question, and questions under topics and subtopics. The survey
provides no topics of its own, so we obtain them by clustering the question
texts.

\paragraph{Embedding.} We embed every node in the Poincar\'e ball
$\mathbb{B}^d_c = \{x \in \mathbb{R}^d : c\lVert x\rVert^2 < 1\}$, the model of
hyperbolic space with curvature $-c$. Hyperbolic space suits trees because
volume grows exponentially with radius, as the number of nodes grows with depth,
so a hierarchy embeds without crowding its lower levels. The ball carries the
geodesic distance
\begin{equation}
  d_c(x, y) = \tfrac{2}{\sqrt{c}}\,
  \operatorname{artanh}\!\bigl(\sqrt{c}\,\lVert (-x) \oplus_c y \rVert\bigr),
\end{equation}
where $\oplus_c$ is M\"obius addition, the ball's analogue of translation.

A tree-recurrent encoder $\phi$ maps each node to the ball, building $\phi(v)$
from $v$'s own text features and the embeddings of its children. Each edge type
$\rho$ (parent--child, axis membership, and so on) has a learned vector
$r_\rho \in \mathbb{B}^d_c$, and a triple $(s, \rho, o)$ is scored by translating
the subject by the relation and measuring distance to the object,
\begin{equation}
  f(s, \rho, o) = -\,d_c\bigl(\phi(s) \oplus_c r_\rho,\ \phi(o)\bigr),
\end{equation}
so a higher score means a shorter edge. We train $\phi$ and the $r_\rho$ by a
margin ranking objective over the tree's edge set $E$: for each true edge we draw
$K$ corrupted negatives $\mathcal{N}(s,\rho,o)$ by replacing an endpoint, and
minimise
\begin{equation}
  \mathcal{L} = \!\!\sum_{(s,\rho,o)\in E}\ \sum_{(s',\rho,o')\in \mathcal{N}}
  \!\!\bigl[\, \gamma + f(s',\rho,o') - f(s,\rho,o) \,\bigr]_+ ,
\end{equation}
with margin $\gamma$ and $[\cdot]_+ = \max(0, \cdot)$. The objective pulls each
node within $\gamma$ of its true neighbours in $d_c$ and pushes corrupted pairs
apart, so related subgroups come to lie close on the ball. We verify the fit
with structural measures: whether depth tracks radius, and tree distance tracks
geodesic distance.

\subsection{Retrieval}

Retrieval looks for places in the graph where subgroups disagree. A \emph{fork}
$f = (a, A, B)$ is an anchor node $a$ with two children $A, B$ whose subtrees
answer differently. It is scored
\begin{equation}
  s(f) = \max\bigl(\mathrm{rel}(q, a),\, 0\bigr)^{\alpha}\; W(A, B).
\end{equation}

\emph{Relevance} $\mathrm{rel}(q, a) = \cos\!\bigl(\psi(q), \psi(a)\bigr)$ is the
cosine similarity between sentence embeddings $\psi$ of the question and the
anchor text, and anchors with $\mathrm{rel}(q,a) < \tau$ are discarded.
Relevance is measured in text space rather than on the ball: the question is not
a graph node and so has no hyperbolic coordinate.

\emph{Divergence} $W(A, B)$ is a Wasserstein distance between the two branches,
each taken as the cloud of its subtree nodes. Let $\{x_i\}_{i=1}^{n}$ and
$\{y_j\}_{j=1}^{m}$ be the nodes under $A$ and $B$, embedded on the ball as
$\phi(x_i), \phi(y_j)$, and give them question-conditioned masses
\begin{equation}
  a_i = \frac{e^{\mathrm{rel}(q, x_i)/\beta}}{\sum_{i'} e^{\mathrm{rel}(q, x_{i'})/\beta}},
  \qquad
  b_j = \frac{e^{\mathrm{rel}(q, y_j)/\beta}}{\sum_{j'} e^{\mathrm{rel}(q, y_{j'})/\beta}},
\end{equation}
so on-topic nodes carry the weight. Then
\begin{equation}
  W(A, B) = \min_{\pi \in \Pi(a, b)}\ \sum_{i=1}^{n} \sum_{j=1}^{m}
            \pi_{ij}\, d_c\bigl(\phi(x_i), \phi(y_j)\bigr),
\end{equation}
where $\Pi(a, b) = \{\pi \ge 0 : \pi\mathbf{1} = a,\ \pi^\top\mathbf{1} = b\}$ is
the set of transport plans with those marginals. The ground cost is the geodesic
$d_c$, so $W$ is the least mass-weighted geodesic displacement needed to move one
branch's opinion cloud onto the other's: large when the branches sit far apart on
the learned manifold. Relevance and divergence thus use different geometries by
design---text-space similarity to judge topicality, hyperbolic distance to judge
disagreement.

A high score marks a relevant anchor whose branches genuinely diverge. Retrieval
returns the surviving forks in decreasing $s$, and both methods below use it
unchanged.

\subsection{Overview}

Both methods write three drafts and merge them into one answer. The first draft
is always a plain answer with no retrieved content; the methods differ only in
how the other two drafts are built, described next. Sampling and merging are
shared, and described in Section~\ref{sec:merge}. As a control, we also run the
same pipeline with three plain drafts and no retrieved content at all, which we
call \emph{plain merge}.

\subsection{Instruction-varied drafts}

In the first method (\textsc{merge\_v2}), drafts~2 and~3 each receive
\emph{every} retrieved fork, and differ in how the forks are rendered. Draft~2
sees each fork as its two contrasting subgroups with their answer
distributions; draft~3 sees each fork's full spectrum of subgroups. Both are
told that the retrieved context is incomplete and may be irrelevant, to judge
which perspectives bear on the question, and to cover the full range of
positions people hold, including ones absent from the context, without
averaging disagreement into a consensus.

Each draft therefore holds all retrieved viewpoints at once, and the model
decides which to foreground and how to balance them. Variation between drafts
comes from rendering and sampling, not from assigning them different content.

\subsection{Subgroup-conditioned drafts}

In the second method (\textsc{persona\_merge}), drafts~2 and~3 each receive
exactly one subgroup, chosen by a fixed rule from the top-ranked fork. Only
that fork is used.

\paragraph{Opinion axis.} Given the top fork $f = (a, A, B)$, let $L(a)$ be the
children of $a$ that carry distributions. Define the pole-to-pole direction
$v = p_B - p_A$ and the scalar coordinate
\begin{equation}
  t(\ell) \;=\; \frac{\langle\, p_\ell - p_A,\; v \,\rangle}
                     {\lVert v \rVert^{2}},
\end{equation}
so that $t(A) = 0$ and $t(B) = 1$. Sorting $L(a)$ by $t$ places the subgroups on
an opinion axis. A subgroup may fall outside the segment, with $t < 0$ or
$t > 1$, so the extremes of the sorted list need not be $A$ and $B$.

\paragraph{Persona selection.} Personas are selected by rule, not by the model:
we take the two endpoints of the sorted axis, and for larger budgets repeatedly
split the widest remaining gap. Endpoints come first because they lie furthest
apart and are likeliest to make the drafts disagree, which is the only thing a
merge can harvest. With two persona drafts the selection is exactly the two
endpoints.

We treat the fork as prior knowledge supplied to the model. Left alone, a model
answers from its own implicit sense of which views exist, and that sense favours
the majority. The fork instead supplies an empirical prior over viewpoints: the
subgroups $L(a)$, their measured distributions $p_\ell$, and their positions on
the axis. Each persona draft is then conditioned on one support point
$\ell_k$ of this prior,
\begin{equation}
  y_k \sim P_M\bigl(\,\cdot \mid q,\ \ell_k\bigr),
  \qquad \ell_k \in L(a),
\end{equation}
where $P_M$ is the model's generation distribution. Selection depends only on a
subgroup's position on the axis, not on how many people share its view, so a
minority position can receive its own draft as readily as a majority one.

\paragraph{Draft prompt.} Each persona draft receives its subgroup rendered as
the group's label and real answer percentages, with options below $5\%$
suppressed, followed by the question. The instruction asks the model to argue
the question as that group would, to commit to the view without hedging, to
present no other group's position, and not to mention the group, the survey,
or the instruction. A draft sees its assigned subgroup and nothing else: not
the fork list, not the ranking, and not the other draft's subgroup.

The anchor is a survey item related to the question, not the question itself,
so a persona is the view of a group that answered a neighbouring item in a
particular way. This is a source of topic drift that the first method shares,
since it retrieves the same forks.

\subsection{Sampling and merging}
\label{sec:merge}

\paragraph{Sampling.} The three drafts are sampled independently at temperature
$0.7$, so no draft sees another. Draft~1 is a plain answer to the question,
identical to a single-answer baseline; we keep it because the plain answer is
the strongest single condition we measure. In the persona method the draft set
is
\begin{equation}
  y_1 \sim P_M(\,\cdot \mid q\,), \qquad
  y_k \sim P_M(\,\cdot \mid q,\ \ell_k\,), \quad k = 2, 3,
\end{equation}
a plain draft alongside two drafts conditioned on opposite ends of the prior.
Each draft covers some viewpoints well, and because they are sampled apart, they
tend to cover different ones.

\paragraph{Merging.} A single further call assembles the drafts into one answer.
The model is asked to act as an editor, not an author: one paragraph per
distinct position, each carried at the length of the draft that explained it
best, and disagreeing positions kept apart rather than reconciled into a middle
view. The aim is to keep the union of what the drafts cover.

Models tend to summarise when asked to combine texts, and summarising loses
exactly the positions a merge is meant to keep. We therefore check the output.
Writing $W_m, P_m$ for the merged answer's words and positions (paragraphs of at
least $20$ words), and $W_d, P_d$ for the largest values over the drafts, the
merge is accepted only if $W_m \ge W_d$ and $P_m \ge P_d$. Otherwise the drafts
are concatenated as they stand. Each method thus costs four model calls per
answer: three drafts and one merge.

\paragraph{Degradation.} When the top anchor carries fewer than three subgroups
with distributions, no axis can be built, no personas are selected, and the
persona method emits the plain draft alone. The instruction-varied method
degrades the same way when retrieval returns no forks. The number of drafts
actually used is recorded per answer, so the resulting dilution toward the
baseline can be measured.

\section{Experimental design}

We run four conditions (Table~\ref{tab:conditions}) and make three comparisons.

\begin{table}[t]
\centering
\small
\begin{tabular}{lll}
\hline
Condition & Drafts 2--3 see & Calls \\
\hline
baseline                & ---                 & 1 \\
plain merge             & nothing             & 4 \\
\textsc{merge\_v2}      & all retrieved forks & 4 \\
\textsc{persona\_merge} & one subgroup each   & 4 \\
\hline
\end{tabular}
\caption{Conditions. The last three share the plain first draft, the sampling,
the merge, the guard, and the call count; they differ only in what drafts~2
and~3 see.}
\label{tab:conditions}
\end{table}

\paragraph{Does drafting and merging help?} Plain merge against the baseline.
Plain merge samples the model three times and merges the samples, with no
retrieved data. Independent samples of one model cover noticeably different
viewpoints, so this step may help on its own, and it has to be measured before
any gain is credited to the survey data.

\paragraph{Does the survey data help?} Each method against plain merge. The two
are matched on everything except what drafts~2 and~3 see: nothing, or retrieved
survey data delivered in that method's way. The difference is therefore the
contribution of the survey data, given the same drafting and merging.

We do not compare the two methods with each other. They differ in several
respects at once (how many retrieved forks they use, how much of that data each
draft sees, and whether the model or a fixed rule picks each draft's viewpoint),
so a direct comparison would not isolate any one of them. Each is instead read
against the same no-data control.

\paragraph{Scoring.} Each question $q$ comes with its surveyed population
partitioned into viewpoint clusters $\mathcal{C}_q$. A judge predicts the rating
$r(u, y) \in \{1, \dots, 5\}$ that participant $u$ would give answer $y$, and a
cluster is \emph{properly represented} when its members endorse the answer on
average:
\begin{equation}
  \mathrm{cov}(c, y) = \mathbf{1}\!\left[\,
    \tfrac{1}{|c|}\textstyle\sum_{u \in c} r(u, y) \ge \theta \,\right],
  \qquad \theta = 4 .
\end{equation}
The score of an answer is the fraction of the question's clusters it represents,
\begin{equation}
  S_w(y, q) = \frac{\sum_{c \in \mathcal{C}_q} w_c\, \mathrm{cov}(c, y)}
                   {\sum_{c \in \mathcal{C}_q} w_c},
\end{equation}
and a condition is scored by $S_w$ averaged over questions. The criterion is
endorsement by a cluster's own members, not similarity to a canonical statement
of its view: naming a position does not represent the people who hold it. We
take $w_c \equiv 1$, the benchmark's own weighting, as the primary metric, and
report $w_c \propto |c|^{-1}$ and a minority restriction as declared-exploratory
secondaries, since uniform weighting counts a view held by a twentieth of the
population exactly as much as one held by a majority.

\paragraph{Judge.} The judge stands in for the surveyed population, so we use
the one whose predicted ratings behave most like human ones. Humans' ratings
barely vary with answer length and endorse about $65\%$ of answers; of two
candidate judges, one rated longer answers increasingly favourably and endorsed
about $75\%$, while the other showed almost no length preference and endorsed
about $65\%$. We use the latter, and flag any result that depends on this
choice.

\paragraph{Reporting.} A coverage difference is a difference of two independent
quantities: clusters gained relative to the reference and clusters lost. We
report both, together with the mean prevalence of each set, because conditions
that tie on the net score can move several times as much mass in each direction.
We also report answer length per condition, since the two methods produce
answers of different lengths and the judge is length-sensitive.

\paragraph{Inference.} All comparisons are paired by question. We bootstrap over
questions, not clusters, since clusters within a question share one answer and
one difficulty; we report percentile intervals, and treat the associated
$p$-values as an inversion of those intervals rather than as independent
evidence. Every condition is judged in a single pass, since participant
subsampling makes scores from separate passes incomparable.

% Bibliography entries for the entire Anthology, followed by custom entries
%\bibliography{custom,anthology-overleaf-1,anthology-overleaf-2}

% Custom bibliography entries only
\bibliography{custom}

\appendix

\section{Example Appendix}
\label{sec:appendix}

This is an appendix.

\end{document}
